%! AP Statistics — Unit 9, Lesson 9.6 — Guided Notes
\documentclass[10pt]{article}
\usepackage{apstats-article}
\usepackage{apstats-boxes}

\begin{document}

\pageheader{Unit 9, Lesson 9.6}{Guided Notes}
\namedateperiod

\vspace{0.1in}

\begin{objectivebox}
By the end of this lesson, I will be able to:
\begin{itemize}
    \item[$\square$] Identify whether a context calls for a confidence interval or a significance test.
    \item[$\square$] Select the correct inference procedure for a given context from Units 6--9.
    \item[$\square$] State $H_0$ and $H_a$ (or describe the parameter being estimated) for the selected procedure.
\end{itemize}
\end{objectivebox}

\begin{vocabbox}
\termblanklong{Parameter}
\termblanklong{Confidence interval}
\termblanklong{Significance test}
\termblanklong{Inference procedure}
\end{vocabbox}

\begin{hookbox}
You have now learned more than 8 different inference procedures. How do you know which one to use on the AP exam? The key is a systematic decision framework.

\textit{Think about it:} What is the first question you should ask when you read a scenario on the AP exam? \writelines{1}
\end{hookbox}

\begin{notesbox}{Decision Framework: Step-by-Step}
Use the three questions below to identify the correct procedure.

\vspace{0.1in}
\textbf{Step 1 — What TYPE of variable(s)?}

\hspace{1cm} $\to$ \blank{4cm} (categories, counts, yes/no) \quad OR \quad \blank{4cm} (measurements, averages)

\vspace{0.15in}
\textbf{Step 2 — What is the DATA STRUCTURE?}

\renewcommand{\arraystretch}{1.5}
\begin{tabular}{@{\hspace{0.5cm}} l l}
  \textit{If Categorical:} & \\
  \hspace{0.5cm} One proportion & $\to$ \blank{5cm} \\
  \hspace{0.5cm} Two proportions (independent) & $\to$ \blank{5cm} \\
  \hspace{0.5cm} One categorical variable, testing fit to a distribution & $\to$ \blank{5cm} \\
  \hspace{0.5cm} Two categorical variables (association / independence) & $\to$ \blank{5cm} \\
  \hspace{0.5cm} One categorical variable compared across 2+ populations & $\to$ \blank{5cm} \\[6pt]
  \textit{If Quantitative:} & \\
  \hspace{0.5cm} One mean & $\to$ \blank{5cm} \\
  \hspace{0.5cm} Two paired means (same subjects, measured twice) & $\to$ \blank{5cm} \\
  \hspace{0.5cm} Two independent means & $\to$ \blank{5cm} \\
  \hspace{0.5cm} Two quantitative variables (predict $y$ from $x$) & $\to$ \blank{5cm} \\
\end{tabular}

\vspace{0.15in}
\textbf{Step 3 — CI or Significance Test?}

\hspace{1cm} \textbf{Confidence interval} $\to$ goal is to \blank{5cm} a parameter.

\hspace{1cm} \textbf{Significance test} $\to$ goal is to \blank{5cm} a claim about a parameter.
\end{notesbox}

\begin{notesbox}{All Inference Procedures: Reference Table}
\renewcommand{\arraystretch}{1.6}
\begin{tabularx}{\linewidth}{@{} p{0.22\linewidth} p{0.25\linewidth} X X @{}}
  \rowcolor{sky}
  \textbf{Variable Type} & \textbf{Data Structure} & \textbf{Confidence Interval} & \textbf{Significance Test} \\
  \hline
  Categorical & 1 proportion & 1-prop $z$-interval & 1-prop $z$-test \\
  Categorical & 2 proportions & 2-prop $z$-interval & 2-prop $z$-test \\
  Categorical & 1 var, $\ge$3 categories (fit to claimed dist.) & --- & Chi-square GOF test \\
  Categorical & 1 var, 2+ independent groups & --- & Chi-square test for homogeneity \\
  Categorical & 2 categorical vars, 1 sample & --- & Chi-square test for independence \\
  Quantitative & 1 mean & 1-sample $t$-interval & 1-sample $t$-test \\
  Quantitative & 2 independent means & 2-sample $t$-interval & 2-sample $t$-test \\
  Quantitative & Matched pairs & Paired $t$-interval & Paired $t$-test \\
  Quantitative & Bivariate (slope) & $t$-interval for slope & $t$-test for slope \\
\end{tabularx}
\end{notesbox}

\begin{notesbox}{Quick Practice}
For each scenario, name the inference procedure.

\vspace{0.1in}
\textbf{a.} A researcher randomly surveys 150 voters and asks them to rate their satisfaction on a scale of 1--10. She wants a 95\% CI for the mean satisfaction score.

\hspace{0.5cm} Procedure: \blank{9cm}

\vspace{0.15in}
\textbf{b.} A study randomly assigns cats to one of two diets and records each cat's weight change (kg) over 6 months. Researchers test whether the diets produce different mean weight changes.

\hspace{0.5cm} Procedure: \blank{9cm}
\end{notesbox}

\end{document}
